\section{Sensado y estimación de estado}
\label{sec:estimacion}

\Cref{th:mapa} resuelve la mitad cartográfica del problema de percepción: las
rejillas en \emph{log-odds} son aditivas y por tanto promediables por consenso.
La otra mitad —saber dónde está cada AMR— quedó como \cref{s:pose}. Esta sección
la desarrolla, y muestra que las dos mitades no son simétricas.

\subsection{Tres sensores y lo que cada uno no puede}

Se distinguen sensores \emph{propioceptivos}, que miden el estado interno, y
\emph{exteroceptivos}, que miden el entorno; y \emph{pasivos} o \emph{activos}
según emitan energía \cite{siegwart2011introduction}.

\begin{table}[!ht]
\centering\small
\caption{Los tres sensores del AMR y su limitación estructural.}
\label{tab:sensores}
\begin{tabularx}{\textwidth}{@{}p{2.7cm}p{2.4cm}p{3.0cm}X@{}}
\toprule
Sensor & Clase & Mide & Lo que no puede \\
\midrule
Encoder de rueda & propioceptivo, pasivo & giro de rueda, luego
  $\Delta s_R,\Delta s_L$ & detectar deslizamiento: mide la rueda, no el suelo \\
Giróscopo (IMU) & propioceptivo, pasivo & velocidad angular $\dot\theta$
  & dar rumbo absoluto: integra, y la deriva del sesgo crece con $t$ \\
Telémetro láser & exteroceptivo, activo & tiempo de vuelo, luego
  alcance y azimut & funcionar sin rasgos: en un pasillo liso no informa a lo
  largo \\
\bottomrule
\end{tabularx}
\end{table}

\begin{observacion}[Las limitaciones son complementarias, no redundantes]
\label{obs:complementarios}
El encoder falla exactamente cuando la rueda desliza, que por
\cref{obs:multiplicador-friccion} es cuando se satura la adherencia —es decir,
durante el empuje intenso, justo cuando la coalición más necesita saber dónde
está—. El láser no falla ahí, pero falla en el pasillo largo y liso donde el
encoder acierta. Fusionar no es una mejora incremental: cada sensor cubre el
modo de fallo del otro.
\end{observacion}

\subsection{Odometría y la ley de propagación del error}

Con $b$ la vía y $\Delta s=(\Delta s_R+\Delta s_L)/2$,
$\Delta\theta=(\Delta s_R-\Delta s_L)/b$, la integración de la pose
$p=(x,y,\theta)^{\!\top}$ es
\begin{equation}
  p^{+}=p+\begin{pmatrix}
    \Delta s\cos(\theta+\Delta\theta/2)\\
    \Delta s\sin(\theta+\Delta\theta/2)\\
    \Delta\theta\end{pmatrix}
  \;=:\;f(p,\Delta).
  \label{eq:odometria}
\end{equation}
El modelo de error supone que los errores de las dos ruedas son independientes y
que su varianza es proporcional al valor absoluto de la distancia recorrida
\cite{siegwart2011introduction}:
\begin{equation}
  \Sigma_\Delta=\operatorname{diag}\bigl(k_R|\Delta s_R|,\ k_L|\Delta s_L|\bigr),
  \label{eq:sigma-delta}
\end{equation}
con $k_R,k_L$ constantes que dependen del robot y del suelo y que se determinan
experimentalmente. La covarianza de la pose se propaga por la ley de primer
orden
\begin{equation}
  \Sigma^{+}=F_p\,\Sigma\,F_p^{\!\top}+F_\Delta\,\Sigma_\Delta\,F_\Delta^{\!\top},
  \qquad
  F_p=\frac{\partial f}{\partial p},\quad F_\Delta=\frac{\partial f}{\partial\Delta}.
  \label{eq:propagacion}
\end{equation}

\begin{proposicion}[La incertidumbre transversal crece con el cubo de la distancia]
\label{pr:cubo}
Bajo \eqref{eq:sigma-delta}, en marcha recta de longitud $d$ la varianza del
rumbo y la varianza longitudinal crecen como $\Theta(d)$, mientras que la
varianza \emph{transversal} crece como $\Theta(d^{3})$.
\end{proposicion}

\begin{proof}
El rumbo es suma de incrementos independientes de varianza proporcional al paso,
luego $\operatorname{var}(\theta)\propto d$. El desplazamiento transversal es
$\int_0^d\theta(u)\,du$, integral de un paseo aleatorio, y su varianza es
$c\int_0^d\!\!\int_0^d\min(u,v)\,du\,dv=\Theta(d^{3})$.
\end{proof}

\begin{observacion}[Comprobación y consecuencia]
\label{obs:cubo-num}
Con $k_R=k_L=10^{-4}$ y $b=0{,}5$~m, una simulación de $4000$ realizaciones da,
al doblar la distancia, razones de varianza de $2{,}0$ en rumbo y de
$7{,}6$--$8{,}6$ en transversal, frente a los valores teóricos $2$ y $8$. A
$16$~m la desviación típica transversal alcanza $1{,}05$~m.

La consecuencia es la razón de ser del telémetro: \emph{ninguna calibración de
$k_R,k_L$ salva la odometría}, porque el problema no es la constante sino el
exponente. Reducir $k$ a la mitad retrasa el instante en que la incertidumbre
alcanza $d_{\min}$ en un factor $2^{1/3}\approx1{,}26$ solamente. Únicamente una
medida exteroceptiva rompe el crecimiento.
\end{observacion}

\subsection{El filtro extendido de Kalman en cinco pasos}

La corrección exteroceptiva se organiza en el ciclo clásico
\cite{siegwart2011introduction}, con $M_i$ los rasgos del mapa:

\begin{enumerate}[leftmargin=2.1em,topsep=0.2em,itemsep=0.25em]
  \item \textbf{Predicción} con \eqref{eq:odometria} y \eqref{eq:propagacion}:
  se obtienen $\hat p^{-}$ y $\Sigma^{-}$.
  \item \textbf{Observación}: del barrido láser se extraen rasgos $z_j$ con
  covarianza $R_j$, obtenida propagando \eqref{eq:propagacion} sobre el ajuste
  de recta.
  \item \textbf{Predicción de medida}: $\hat z_i=h(\hat p^{-},M_i)$.
  \item \textbf{Emparejamiento}: innovación $\nu_{ij}=z_j-\hat z_i$ con
  covarianza $\Sigma_{\mathrm{IN}}=H\Sigma^{-}H^{\!\top}+R_j$; se acepta si
  $\nu^{\!\top}\Sigma_{\mathrm{IN}}^{-1}\nu\le\chi^2_{\alpha}$.
  \item \textbf{Estimación}: con $K=\Sigma^{-}H^{\!\top}\Sigma_{\mathrm{IN}}^{-1}$,
  $\hat p=\hat p^{-}+K\nu$ y $\Sigma=(I-KH)\Sigma^{-}$.
\end{enumerate}

\begin{observacion}[El paso frágil es el cuarto]
\label{obs:emparejamiento}
Los pasos 1, 3 y 5 son álgebra. El paso 4 es una decisión discreta, y un
emparejamiento erróneo no degrada la estimación: la corrompe, porque inyecta una
innovación grande con covarianza pequeña. En un almacén —donde las estanterías
son deliberadamente idénticas y periódicas— la ambigüedad de asociación es la
norma y no la excepción. De ahí que el umbral $\chi^2$ deba fijarse conservador
y que convenga rechazar el ciclo completo antes que aceptar un emparejamiento
dudoso.
\end{observacion}

\subsection{Localización cooperativa y SLAM distribuido}

Cuando dos AMR se miden mutuamente, la medida relativa informa a ambos. Fusionar
promediando covarianzas es incorrecto: tras el primer intercambio las
estimaciones dejan de ser independientes, y el promedio vuelve a contar la
información ya compartida, produciendo una covarianza optimista.

\begin{proposicion}[Fusión consistente sin conocer la correlación]
\label{pr:interseccion}
Sean $(\hat p_1,\Sigma_1)$ y $(\hat p_2,\Sigma_2)$ dos estimaciones de la misma
pose con correlación cruzada desconocida. La intersección de covarianzas
\[
  \Sigma_{\mathrm{CI}}^{-1}=\omega\,\Sigma_1^{-1}+(1-\omega)\,\Sigma_2^{-1},
  \qquad
  \Sigma_{\mathrm{CI}}^{-1}\hat p_{\mathrm{CI}}
  =\omega\,\Sigma_1^{-1}\hat p_1+(1-\omega)\,\Sigma_2^{-1}\hat p_2,
\]
con $\omega\in[0,1]$, es consistente para \emph{cualquier} correlación
\cite{julier1997ci}. El precio es que $\Sigma_{\mathrm{CI}}$ domina a la
covarianza que se obtendría conociendo la correlación: se paga conservadurismo a
cambio de no necesitar el término cruzado.
\end{proposicion}

\begin{teorema}[Asimetría entre cartografiar y localizar]
\label{th:asimetria}
La cartografía cooperativa y la localización cooperativa no admiten el mismo
tratamiento distribuido:
\begin{enumerate}[label=(\roman*),leftmargin=2.2em,topsep=0.2em,itemsep=0.15em]
  \item el mapa \emph{sí} consensúa: por \cref{le:aditividad} la evidencia en
  \emph{log-odds} es aditiva, luego el promedio de consenso de \cref{th:mapa}
  converge a la fusión bayesiana exacta;
  \item la pose \emph{no}: las poses viven en $SE(2)$, donde el promedio
  aritmético no es una operación intrínseca, y además las estimaciones se
  correlacionan al intercambiarse, de modo que el promedio subestima la
  covarianza.
\end{enumerate}
En consecuencia la cartografía admite un protocolo de consenso con garantía,
mientras que la localización cooperativa admite solo fusión conservadora
(\cref{pr:interseccion}) o estimación conjunta centralizada.
\end{teorema}

\begin{proof}
(i) es \cref{le:aditividad} junto con \cref{th:mapa}. (ii) tiene dos causas
independientes: la variedad —el promedio de dos rotaciones no es en general una
rotación, y la linealización que lo repara vale solo localmente— y la
correlación, que invalida el promedio incluso en el caso lineal, según
\cref{pr:interseccion}.
\end{proof}

\begin{observacion}[Por qué la asimetría no rompe la arquitectura]
\label{obs:asimetria-consecuencia}
\Cref{th:asimetria} no obliga a centralizar, porque las dos mitades entran en el
mecanismo por vías distintas. El mapa entra como bien común: alimenta el
catálogo de rutas, y ahí el consenso es correcto y demostrable. La pose entra
como dato local: cada AMR necesita \emph{su} pose para evaluar su barrera, y de
los demás no necesita la pose global con precisión métrica, sino la distancia
relativa, que el láser mide directamente sin pasar por dos estimaciones
globales. La medida relativa esquiva el problema de fusión justo en el punto
donde importaría.
\end{observacion}

\begin{observacion}[Qué nivel de mapa usa cada capa]
\label{obs:jerarquia-mapas}
La navegación en interiores distingue tres niveles de representación
\cite{gorrostieta2019applications}: \emph{geométrico} —métrico, la rejilla—,
\emph{topológico} —nodos y adyacencias— y \emph{semántico} —etiquetas y
relaciones sobre las que razonar—. Este trabajo usa los dos primeros y no el
tercero: la rejilla de \cref{th:mapa} es geométrica y alimenta las barreras;
el catálogo de rutas y el grafo de arcos de \cref{th:multiclase} son
topológicos y alimentan el peaje.

La ausencia del nivel semántico es deliberada y tiene consecuencia. Sin él, la
prioridad de una carga entra como el escalar de urgencia del pago y no como una
propiedad del objeto —«frágil», «refrigerado», «no inclinable»—, y restricciones
de ese tipo tendrían que traducirse a mano a cotas sobre el \emph{twist}
admisible. Es una vía de extensión natural, no un fundamento que falte.
\end{observacion}

\subsection{De la covarianza al margen de seguridad}

\Cref{pr:pose} tensa las barreras por la incertidumbre de pose sin decir cuánto.
El estimador lo cuantifica.

\begin{corolario}[Tensado derivado del filtro]
\label{co:tensado}
Si el filtro entrega $\Sigma$ y se exige que la barrera se respete con
probabilidad $1-\alpha$ bajo aproximación gaussiana, basta tensar la holgura en
\[
  d_{\mathrm{tens}}=\sqrt{\chi^2_{2,\,1-\alpha}\;\lambda_{\max}\bigl(\Sigma_{xy}\bigr)},
\]
con $\Sigma_{xy}$ el bloque posicional. Sustituyendo esta cantidad en
\cref{le:astop} y en \cref{th:astop-cuerpo}, las garantías de invariancia se
mantienen evaluadas sobre la pose estimada.
\end{corolario}

\begin{observacion}[Lo que el tensado cuesta y lo que cierra]
\label{obs:tensado-coste}
El tensado no es gratis: consume holgura, y por \cref{obs:precio-astop} la
holgura fija la velocidad. Un filtro que degrada —porque el láser no ve rasgos o
porque el emparejamiento se rechaza— eleva $\lambda_{\max}(\Sigma)$, tensa más y
frena la flota. Esa cadena convierte la calidad de la estimación en un
\emph{recurso} más, con la misma estructura que $(a,\delta,\eta)$ en
\cref{th:degradacion}: la incertidumbre no provoca un fallo, provoca degradación
continua del caudal.

Con esto \cref{s:pose} deja de ser una hipótesis en blanco: se declara qué
estimador la sostiene, cuánta incertidumbre produce y cómo esa incertidumbre
entra en las garantías. Lo que sigue supuesto —y así se declara— es la
disponibilidad de rasgos suficientes para que el paso 4 no se rechace de forma
sostenida.
\end{observacion}

\begin{observacion}[Estimador y mecanismo se han diseñado por separado]
\label{obs:observador-mecanismo}
Queda una limitación estructural que conviene nombrar. Este trabajo diseña el
estimador en esta sección y el mecanismo de precios en las anteriores, cada uno
por su cuenta, y después acopla ambos por la vía del tensado de
\cref{co:tensado}. Existe una formulación que los trata como un solo problema
—\emph{diseño conjunto de observador y mecanismo}— cuya motivación es
precisamente que el rendimiento del mecanismo depende de lo que el observador
puede ver, y que por tanto elegirlos por separado desperdicia grados de libertad
\cite{clempner2024optimization}.

La consecuencia práctica en este sistema sería, por ejemplo, sesgar el
emparejamiento del paso 4 hacia los rasgos que más reducen la incertidumbre
\emph{en las direcciones que la barrera tensa}, en lugar de hacia los que más
reducen la traza de la covarianza. No es lo mismo: la barrera solo consume la
holgura en la dirección de aproximación al obstáculo. Se declara como diseño
separado, con la vía de mejora identificada y no explorada.
\end{observacion}
